\documentclass[english,compress]{beamer}
\usepackage{kloeckislides}
\nonstopmode

\usepackage[normalem]{ulem}
\usepackage{pifont}
\usepackage{ifthen}

\setbeamercolor{section in head/foot}{use=structure,bg=structure.fg!25!bg}
\defbeamertemplate*{footline}{split theme}
{%
  \leavevmode%
  \begin{beamercolorbox}[wd=.5\paperwidth,ht=2.5ex,dp=1.125ex]{section in head/foot}%
    \insertsectionnavigationhorizontal{\paperwidth}{\hskip0pt plus1filll}{}%
  \end{beamercolorbox}%
  %\begin{beamercolorbox}[wd=.5\paperwidth,ht=2.5ex,dp=1.125ex]{subsection in head/foot}%
    %\insertsubsectionnavigationhorizontal{.5\paperwidth}{}{\hskip0pt plus1filll}%
  %\end{beamercolorbox}%
}


%\useoutertheme[subsection=false]{miniframes}

\setbeamertemplate{frametitle}[default][center]

\AtBeginDocument{%
  {
    \usebeamercolor{section in head/foot}
  }
  
  \pgfdeclareverticalshading{beamer@headfade}{\paperwidth}
  {%
    color(0cm)=(bg);
    color(1.25cm)=(section in head/foot.bg)%
  }

  \setbeamercolor{section in head/foot}{bg=}
}

\addtoheadtemplate{\pgfuseshading{beamer@headfade}\vskip-1.25cm}{}

\beamertemplatedotitem

\setbeamercolor{section in head/foot}{parent=palette quaternary}
\setbeamercolor{subsection in head/foot}{parent=palette primary}

\setbeamercolor{author in head/foot}{parent=section in head/foot}
\setbeamercolor{title in head/foot}{parent=subsection in head/foot}



\AtBeginSection[] {
  \begin{frame}<beamer>
  \frametitle{Outline}
  \tableofcontents[sectionstyle=show/shaded,subsectionstyle=show/show/hide]
\end{frame}
}
\AtBeginSubsection[] {
  \begin{frame}<beamer>
  \frametitle{Outline}
  \tableofcontents[sectionstyle=show/shaded,subsectionstyle=show/shaded/hide]
\end{frame}
}

\newcommand{\technicality}[2]{%
  {\strut #1\\
    \begin{beamercolorbox}[sep=1mm]{block body}
      #2
    \end{beamercolorbox}
  }%
}

\lstset{
  language=C++,
  rangebeginprefix=//\ ,
  rangeendprefix=//\ ,
}

\def\weblink#1#2{\href{#1}{\color{blue}\underline{#2}}}

\definecolor{fetch}{RGB}{227,110,35}
\definecolor{alu}{RGB}{255,188,24}
\definecolor{context}{RGB}{132,146,175}

\usepackage{keystroke}

\setbeamertemplate{navigation symbols}{}


\begin{document}
% {{{ front matter

\title{High-Performance Scientific Computing\\Lecture 4: OpenCL}

\date{MATH-GA 2011 / CSCI-GA 2945 $\cdot$ September 26, 2012}

\frame{\titlepage}

\begin{frame}{Today}
  \tableofcontents[hideallsubsections]
\end{frame}
% }}}
% -----------------------------------------------------------------------------
\begin{frame}{Admin Bits}
  \begin{itemize}[<+->]
    \item HW1 graded before weekend
    \item HW2 due
    \item HW3 out
  \end{itemize}
\end{frame}
% -----------------------------------------------------------------------------
\section{HW2}
% -----------------------------------------------------------------------------
\begin{frame}{HW2 problem 2}
  \begin{center}
  \Huge Demo time
  \end{center}
\end{frame}
% -----------------------------------------------------------------------------
\begin{frame}{OpenMP sync primitives}
  \begin{itemize}[<+->]
    \item Critical section
    \item Locks
    \item Atomics
      \begin{itemize}
        \item Update: \texttt{x++;}
        \item Capture: \texttt{v = x++;}
        \item Structured: {v = x; x |= expr;} (``Test-and-set'')
        \item Compare-and-swap (not in OpenMP)
      \end{itemize}
  \end{itemize}
\end{frame}
% -----------------------------------------------------------------------------
\section{Chips for Throughput}
% -----------------------------------------------------------------------------
% {{{
\input{cpu-chip-real-estate}

\newcommand{\kayvoncredit}{
  \begin{tikzpicture}[overlay]
    \node [xshift=1cm,yshift=0.5cm]
      at (current page.south west)
      [font=\scriptsize,fill=gray!30,anchor=south west,opacity=0.5]
      {Credit: Kayvon Fatahalian (Stanford) };
  \end{tikzpicture}
}
\newcommand{\kayvonframe}[5]{

  \begin{frame}{#1}
    #4
    \begin{center}
    \includegraphics[viewport=#3,clip=true,page=#2,height=0.7\textheight]{kayvon-gpuarch.pdf}
    \end{center}
    \kayvoncredit
    #5
  \end{frame}
}

\kayvonframe{``CPU-style'' Cores}{13}{1in 1in 9in 6.5in }{}{}
\kayvonframe{Slimming down}{14}{1in 1in 9in 6.5in }{}{}
\kayvonframe{More Space: Double the Number of Cores}{15}{2.5in 1in 8.5in 6.5in }{}{}
\kayvonframe{\dots again}{16}{4in 1in 9.5in 6.5in }{}{}
\kayvonframe{\dots and again}{17}{4in 1.35in 9.5in 6.5in }{}{
  \uncover<2>{
    \begin{tikzpicture} [overlay]
      \node [above left=1cm of current page.south east, draw,drop shadow,fill=white,
      text width=0.6\textwidth, inner xsep=0.5cm,inner ysep=0.5cm,thick]
        {
          $\rightarrow$ 16 independent instruction streams

          \medskip
          Reality: instruction streams not actually
          very different/independent
        } ;
    \end{tikzpicture}
  }
}
\begin{frame}{Saving Yet More Space}
  \begin{columns}
    \column{0.5\textwidth}
      \only<1-2>{%
      \includegraphics[viewport=1.8in 3.8in 5.45in 6.25in,clip=true,page=19,width=\textwidth]{kayvon-gpuarch.pdf}\\[-0.5mm]
      \includegraphics[viewport=1.8in 1.35in 5.45in 3.8in,clip=true,page=19,width=\textwidth]{kayvon-gpuarch.pdf}
      }%
      \only<3>{%
      \includegraphics[viewport=1.8in 3.8in 5.45in 6.25in,clip=true,page=20,width=\textwidth]{kayvon-gpuarch.pdf}\\[-0.5mm]
      \includegraphics[viewport=1.8in 1.35in 5.45in 3.8in,clip=true,page=19,width=\textwidth]{kayvon-gpuarch.pdf}
      }%
      \only<4->{%
      \includegraphics[viewport=1.8in 3.8in 5.45in 6.25in,clip=true,page=20,width=\textwidth]{kayvon-gpuarch.pdf}\\[-0.5mm]
      \includegraphics[viewport=1.8in 1.35in 5.45in 3.8in,clip=true,page=20,width=\textwidth]{kayvon-gpuarch.pdf}
      }
    \column{0.5\textwidth}%
      \uncover<2->{%
        \textbf{Idea \#2}

        \medskip
        Amortize cost/complexity of managing an instruction
        stream across many ALUs

        \medskip
        \large \textbf{$\rightarrow$ SIMD}
      }
  \end{columns}
  \kayvoncredit
\end{frame}
\kayvonframe{Gratuitous Amounts of Parallelism!}{24}{5in 1.45in 10in 6.75in }{}{
  \uncover<2->{
    \begin{tikzpicture} [overlay]
      \node [below right=1.75cm of current page.north west, draw,drop shadow,fill=white,
      text width=0.8\textwidth, inner sep=2.5mm,thick]
        {
          Example:

          \medskip
          128 instruction streams in parallel

          16 independent groups of 8 synchronized streams
        } ;
    \end{tikzpicture}
  }
  \uncover<3>{
    \begin{tikzpicture} [overlay]
      \node [above left=1cm of current page.south east, draw,drop shadow,fill=white,
      text width=0.6\textwidth, inner xsep=0.5cm,inner ysep=0.5cm,thick]
        {
          Great if everybody in a group does the same thing.

          \medskip
          But what if not?

          \medskip
          What leads to divergent instruction streams?
        } ;
    \end{tikzpicture}
  }
}
\kayvonframe{Branches}{26}{0.85in 0.9in 10.5in 6.8in }{}{}
\kayvonframe{Branches}{27}{0.85in 0.9in 10.5in 6.8in }{}{}
\kayvonframe{Branches}{28}{0.85in 0.9in 10.5in 6.8in }{}{}
\kayvonframe{Branches}{29}{0.85in 0.9in 10.5in 6.8in }{}{}

\begin{frame}{Recent Processor Architecture}
  \begin{columns}
    \column{0.5\textwidth}
      \begin{itemize}
      \item Commodity chips
      \item ``Infinitely'' many cores
      \item ``Infinite'' vector width
      \item Must hide memory latency\\
        ($\rightarrow$ ILP, SMT)
      \end{itemize}
    \column{0.5\textwidth}
      \begin{itemize}
      \item Compute bandwidth\\
        \hfill $\gg$ Memory bandwidth
      \item Bandwidth only achievable by \emph{homogeneity}
      \end{itemize}
      \vspace*{3.5ex}
  \end{columns}
  \hspace*{-0.25\textwidth}
  \begin{tabular}{p{0.25\textwidth}p{0.25\textwidth}p{0.25\textwidth}p{0.25\textwidth}p{0.25\textwidth}}
  \includegraphics[width=0.25\textwidth]{gt200-die.jpg}
  &
  \includegraphics[width=0.25\textwidth]{fermi-die-shot.jpeg}
  &
  \centering
  \includegraphics[width=3.5cm,angle=90]{ivy-bridge-die-shot.jpeg}
  &
  \includegraphics[width=0.25\textwidth]{tahiti-die-shot.jpeg}
  &
  \includegraphics[width=0.25\textwidth]{gk110-die-shot.jpeg}
  \\
  \centering Nv~GT200\par(2008)
  &
  \centering Nv~Fermi\par(2010)
  &
  \centering Intel IVB\par(2012)
  &
  \centering AMD Tahiti \par(2012)
  &
  \centering Nv GK110\par(2012?)
  \end{tabular}
\end{frame}
% }}}
% -----------------------------------------------------------------------------
\section{OpenCL: Overview}
% -----------------------------------------------------------------------------
% {{{
\begin{frame}{What is OpenCL?}
  \begin{columns}
    \column{0.7\textwidth}

      OpenCL (Open Computing Language) is an open, royalty-free
      standard for general purpose parallel programming across CPUs,
      GPUs and other processors.
      \hfill{\footnotesize[OpenCL 1.1 spec]}
      \bigskip


      \begin{itemize}
        \item Device-neutral (Nv GPU, AMD GPU, Intel/AMD CPU)
        \item Vendor-neutral
        \item Comes with `JIT' compilation
      \end{itemize}
      Defines:
      \begin{itemize}
        \item Host-side programming interface (library)
        \item Device-side programming language (!)
      \end{itemize}

    \column{0.3\textwidth}
      \includegraphics[width=\textwidth] {opencl-logo.png}

  \end{columns}
\end{frame}

\newcommand{\khronoscredit}{
  \begin{tikzpicture}[overlay]
    \node [xshift=1cm,yshift=0.5cm]
      at (current page.south west)
      [font=\scriptsize,fill=gray!30,anchor=south west,opacity=0.5]
      {Credit: Khronos Group};
  \end{tikzpicture}
}
\def\khronosslide#1#2
{
  \begin{frame}{#1}
    \hspace*{-0.75cm}\includegraphics[viewport=2cm 0cm 31cm 14.5cm,clip=true,width=1.15\textwidth,page=#2]{opencl-overview.pdf}
    \khronoscredit
  \end{frame}
}

\input{cl-vocabulary}

%\khronosslide{OpenCL: Execution Model}{11}

{
\def\drawgroup#1{
  \foreach \i in {0,1,...,23}
    \foreach \j in {0,1,...,16}
    {
      \pgfmathtruncatemacro{\grp}{(\i+24*\j)/16}
      \ifthenelse{\equal{\grp}{#1}}
        {\def\grpcolor{gray!60}}
        {\def\grpcolor{gray!30}}
      \draw [line width=0.5pt,fill=\grpcolor] (\i,\j) rectangle +(1,1) ;
    }
}

\def\simplesm{{
  \node (fetch) [blk,fill=fetch] at (0,0) {Fetch/ Decode} ;
  \begin{scope}[yshift=-1.25cm]
  \foreach \x in {0,1,2,3}
  \foreach \y in {1,0}
  \draw [fill=alu,thick] (0.625*\x,0.5*\y) rectangle +(0.625,0.5);
  \end{scope}
  \node (pvt) [blk,fill=context, minimum height=2cm,below=1.5cm of fetch]
  {32 kiB Ctx Private (``Registers'') };
  \node (shared) [blk,fill=context,minimum height=1cm,below=0.2cm of pvt]
  {16 kiB Ctx Shared };
}}
\def\drawprocassign#1#2#3{
  \pgfmathtruncatemacro{\gj}{#1/4}
  \pgfmathtruncatemacro{\gi}{#1-\gj*4}
  \pgfmathtruncatemacro{\si}{#2/3}
  \pgfmathtruncatemacro{\sj}{#2-\si*3}
  \draw [->,ultra thick,#3] (g\gi\gj) ++(0.5,-0.5) ..controls +(0.3,1.2-0.2*\gi) and +(-0.5,0.5).. (sm\sj\si) ;
}
\begin{frame}{Connection: Hardware $\leftrightarrow$ Programming Model }
  \hspace*{-3mm}
  \begin{tikzpicture}[
  blk/.style={draw,thick,text centered,anchor=south west,minimum width=2.5cm,text width=2cm,font=\small},
  ]
    \uncover<+->{
      \begin{scope}[xshift=8cm]
        \draw [thick,fill=black!20] (-0.5,0.5) rectangle +(3.5,-5.5)
        coordinate (hwboxse);

        \foreach \i in {0,1,2}
        \foreach \j in {0,1,2}
        {
          \node at (\i,-\j*1.75) (sm\i\j) {
            \begin{tikzpicture}[transform canvas={scale=0.25}]
              \simplesm
            \end{tikzpicture}
          };
        }
      \end{scope}
    }

    \uncover<+-+(2)>{
        \node [rotate=20,font=\Large\bfseries,
        cloud,cloud ignores aspect=true,cloud puffs=15,draw,thick,
        text width=4cm, text centered,anchor=north
        west,xshift=1cm,yshift=-1.25cm]
        {Who cares how many cores? };
    }
    \uncover<+-+(1)>{ 
      \node [anchor=west,draw,drop shadow,fill=white,
      text width=0.5\textwidth, inner xsep=0.5cm,inner ysep=0.5cm,thick]
      (prog-model-idea)
      at (0,-4cm)
        {
          Idea:

          \begin{itemize}
            \item Program as if there were ``infinitely'' many cores
            \item Program as if there were ``infinitely'' many ALUs per core
          \end{itemize}
        } ;
    }
    \uncover<+>{
      \node [anchor=north west,xshift=0.3cm,yshift=-0.6cm, draw,drop shadow,fill=white,
      text width=0.8\textwidth, inner xsep=0.5cm,inner ysep=0.5cm,thick]
      at (prog-model-idea.north west)
        {
          Consider: Which is easy to do automatically?
          \begin{itemize}
            \item Parallel program $\rightarrow$ sequential hardware
          \end{itemize}
          or
          \begin{itemize}
            \item Sequential program $\rightarrow$ parallel hardware?
          \end{itemize}

        } ;
    }
    \uncover<+>{ % notes disappear
    }
    \uncover<+->{
      \draw [thick,|->] (0,0.3) -- ++(3,0) 
        node [anchor=west,font=\small] { Axis 0 };
      \draw [thick,|->] (-0.3,0) -- ++(0,-2) 
        node [anchor=east,font=\small,rotate=90] { Axis 1 };

      \foreach \gi in {0,1,2,3}
        \foreach \gj in {0,1,2}
        {
          \draw [very thick,fill=black!20] (\gi*1.75,\gj*-1.5) coordinate (g\gi\gj)
            rectangle +(1.75,-1.5) ;
          \foreach \li in {1,2,3,4,5,6}
            \draw (g\gi\gj) ++(0.25*\li,0) -- +(0,-1.5) ;
          \foreach \li in {1,2,3,4,5}
            \draw (g\gi\gj) ++(0,-0.25*\li) -- +(1.75, 0) ;
        }
    }
    \coordinate (gridcenter) at (3,-2.25);
    \uncover<+->{
      \node [below=1.5cm of sm12,font=\Large] (hwlabel) { Hardware };
    }
    \uncover<+->{
      \node [below=2.5cm of gridcenter,font=\Large] (swreplabel)
      { Software representation };
    }
    \only<+-+(2)>{
      \node [font=\huge\bfseries,rotate=10, text width=5cm] at (3,-2.25) { Grid 

      (Kernel: Function on Grid) };
    }
    \only<+-+(1)>{
      \fill [black,opacity=0.3] (g00) rectangle (g11) ;
      \node [font=\Large\bfseries,anchor=north west] at (g10) { (Work) Group};
    }
    \only<+>{
      \fill [black,opacity=0.3] (g32) ++(0.5,-0.25) rectangle +(0.25,-0.25)
      node [pos=0.5] (workitem) {};
      \node [draw,thick,fill=white,font=\Large\bfseries,below=0.25 of
        workitem, arrow box, arrow box arrows={north:0.5cm}] { (Work) Item};
    }
    \only<+>{}
    \uncover<+>{
      \draw [<->,line width=1mm] (gridcenter)
      ..controls +(1,1) and +(-1,1)..
      (sm11)
      node [pos=0.6,anchor=south,yshift=0.5cm,font=\fontsize{40}{30}\bfseries] {?};
    }
    \only<+>{
      \fill [black,opacity=0.3] (g00) rectangle (g11) ;
      \drawprocassign{0}{0}{}
    }
    \only<+>{ \foreach \gnum in {0,1,2} { \drawprocassign{\gnum}{\gnum}{} } }
    \only<+>{ \foreach \gnum in {0,1,...,8} { \drawprocassign{\gnum}{\gnum}{} } }
    \only<+>{ \foreach \gnum in {0,1,...,8} { \drawprocassign{\gnum}{\gnum}{opacity=0.2} } }
    \only<+>{
      \foreach \gnum in {0,1,...,8}
      {
        \drawprocassign{\gnum}{\gnum}{opacity=0.2}
        \pgfmathtruncatemacro{\gj}{\gnum/4}
        \pgfmathtruncatemacro{\gi}{\gnum-\gj*4}
        \foreach \id in {0,...,7}
        {
          \pgfmathtruncatemacro{\li}{\id/7}
          \pgfmathtruncatemacro{\lj}{\id-\li*7}
          \fill [black,opacity=0.3]
            (g\gi\gj) ++ (0.25*\lj,-0.25*\li) rectangle +(0.25,-0.25);
        }
      }
    }
    \only<+>{
      \foreach \gnum in {0,1,...,8}
      {
        \drawprocassign{\gnum}{\gnum}{opacity=0.2}
        \pgfmathtruncatemacro{\gj}{\gnum/4}
        \pgfmathtruncatemacro{\gi}{\gnum-\gj*4}
        \foreach \id in {8,...,15}
        {
          \pgfmathtruncatemacro{\li}{\id/7}
          \pgfmathtruncatemacro{\lj}{\id-\li*7}
          \fill [black,opacity=0.3]
            (g\gi\gj) ++ (0.25*\lj,-0.25*\li) rectangle +(0.25,-0.25);
        }
      }
    }
    \only<+-+(2)>{
      \foreach \gnum in {0,1,...,8}
      {
        \drawprocassign{\gnum}{\gnum}{opacity=0.2}
        \pgfmathtruncatemacro{\gj}{\gnum/4}
        \pgfmathtruncatemacro{\gi}{\gnum-\gj*4}
        \foreach \id in {16,...,23}
        {
          \pgfmathtruncatemacro{\li}{\id/7}
          \pgfmathtruncatemacro{\lj}{\id-\li*7}
          \fill [black,opacity=0.3]
            (g\gi\gj) ++ (0.25*\lj,-0.25*\li) rectangle +(0.25,-0.25);
        }
      }
    }
    \uncover<+>{
        \node [below left=1cm of hwboxse, draw,drop shadow,fill=white,
        text width=0.6\textwidth, inner xsep=0.5cm,inner ysep=0.5cm,thick,
        yshift=2cm]
          {
            Really: Group provides pool of parallelism to draw from.

            \medskip
            X,Y,Z order \emph{within} group matters. (Not \emph{among}
            groups, though.)
          } ;
    }
    \only<+>{}
    \only<+>{ \foreach \gnum in {9,10,11} { \drawprocassign{\gnum}{(\gnum-9)}{} } }
    \only<+->{
      \fill [black,opacity=0.3] (g31) ++(+1,-0.25) rectangle ++(0.25,-0.25)
      coordinate [pos=0.5] (workitem2);
      \node [draw,thick,fill=white,below=0.25 of
        workitem2, arrow box, arrow box arrows={north:0.5cm},text width=0.85\textwidth] { 
          \begin{itemize}
            \item \texttt{get\_local\_id(axis)?/size(axis)?}
            \item \texttt{get\_group\_id(axis)?/num\_groups(axis)?}
            \item \texttt{get\_global\_id(axis)?/size(axis)?}
          \end{itemize}
          \texttt{axis=0,1,2,\dots}
        };
    }
  \end{tikzpicture}

  % ---------------------------------------------------------------------------
  \uncover<+>{%
    \begin{tikzpicture} [overlay]
      \node [below right=1cm of current page.north west, draw,drop shadow,fill=white,
      text width=0.6\textwidth, inner xsep=0.5cm,inner ysep=0.5cm,thick]
        {
          Grids can be 1,2,3-dimensional.
        } ;
    \end{tikzpicture}%
  }
\end{frame}




\begin{frame}{Dive into OpenCL: Preparation}
  \begin{center}
  \Huge Demo time
  \end{center}
\end{frame}

% }}}
% -----------------------------------------------------------------------------
\section{OpenCL: Between host and device}
% -----------------------------------------------------------------------------
\input{cl-command-queues}
% -----------------------------------------------------------------------------
\section{OpenCL: Device Language}
% -----------------------------------------------------------------------------
% {{{
\begin{frame}{OpenCL Device Language}
  \begin{columns}
    \column{0.7\textwidth}
      OpenCL device language is C99, with these differences:

      \medskip
      \plusball Index getters

      \plusball Memory space qualifiers

      \plusball Vector data types 

      \plusball Many generic (`overloaded') math functions

      \plusball Synchronization

      \minusball Recursion

      \minusball Fine-grained \texttt{malloc()}

      \minusball Function pointers
    \column{0.3\textwidth}
      \includegraphics[width=\textwidth]{opencl-logo.png}
  \end{columns}
\end{frame}

%\input{cuda-cl-dictionary}

\begin{frame}{Address Space Qualifiers}

  \begin{center}
  \begin{tabular}{p{5em}cccp{2.8cm}}
    \hline
    \textbf{Type} & \textbf{Per} & \textbf{``Speed''} \\
    \hline
    private*) & work item & super-fast\\
    local & group & fast \\
    global & grid & kinda slow \\
    \hline
  \end{tabular}

  \bigskip
  *) default, so optional
  \end{center}
  \uncover<2>{
    \begin{tikzpicture} [overlay]
      \node [
        above left=1cm of current page.south east, draw,drop shadow,fill=white,
        inner sep=5mm,thick, text width=0.6\textwidth]
        {
          Should really discuss ``speed'' in terms of
          latency/bandwidth.

          \bigskip
          \emph{Both} decrease with distance from the point of
          execution.
        } ;
    \end{tikzpicture}
  }
\end{frame}

% }}}
% -----------------------------------------------------------------------------
\section{OpenCL: Synchronization}
% -----------------------------------------------------------------------------
% {{{
\begin{nologo}
\begin{frame}{Concurrency and Synchronization}
  \uncover<+->{
    \begin{beamercolorbox}[sep=3mm]{block body}
      GPUs have layers of concurrency.\\
      \hfill Each layer has its synchronization primitives.
    \end{beamercolorbox}
  }
  \bigskip
  \begin{columns}[c]
    \column{0.6\textwidth}
      \uncover<+->{
        \begin{itemize}
          \item Intra-group:\\ 
            \texttt{barrier(\dots)}, \\
            \texttt{mem\_fence(\dots)}\\
            \texttt{\dots} =
            \texttt{CLK\_\{LOCAL,GLOBAL\}\_MEM\_FENCE}
          \item Inter-group:\\ Kernel launch
          \item CPU-GPU:\\ Command queues, Events
        \end{itemize}
      }
    \column{0.35\textwidth}
      \includegraphics[width=\textwidth]{onion.jpeg}
  \end{columns}

\end{frame}
\end{nologo}
\addimgcredit{Onions: flickr.com/darwinbell \cc}

\input{barrier}
\input{memory-fence}

% -----------------------------------------------------------------------------
\begin{frame}{Synchronization between Groups}
  \begin{block}<+->{Golden Rule:}
    Results of the algorithm must be independent of the order in which
    work groups are executed.
  \end{block}

  \uncover<+->{
  \textbf{Consequences:}
  \begin{itemize}
    \item Work groups may read the same information from global memory.
    \item But: Two work groups may not validly write different things to 
      the same global memory.
    \item Kernel launch serves as
      \begin{itemize}
        \item Global barrier
        \item Global memory fence
      \end{itemize}
  \end{itemize}
  }
\end{frame}
% -----------------------------------------------------------------------------
\input{cl-atomic}
% }}}
\questionframe{}
\imagecreditslide

\end{document}
% vim: foldmethod=marker
